\section{Results}\label{sec:results}

Results are presented in six blocks: uncontrolled drift, perfect-navigation
controller baselines, the estimation-in-the-loop reliability map, the
filter-consistency analysis, the corrected trade curves and knee summary,
and the fidelity-in-the-loop reading. Model validation
(Fig.~\ref{fig:validation}) was presented with the models in
Section~\ref{sec:models}.

\subsection{Uncontrolled drift characterization}

Figures~\ref{fig:drift} and~\ref{fig:contour} quantify the problem the
controllers must solve. Uncontrolled differential drift grows with formation
size and with the differential-element content of the geometry, varies
strongly with chief inclination through the $J_2$ secular rates, and is
dominated by the along-track runaway induced by relative semi-major-axis and
inclination-vector offsets, which is the classical structure
\cite{schaub2001j2invariant,alfriend2010formation}. At 1-km scale the free
drift spans more than two orders of magnitude across geometries. Pure
along-track separation drifts only 0.6--0.9~m/day, tens of metres per month,
whereas the passively safe E/I-vector geometry, which deliberately carries
relative eccentricity and inclination vectors, is the fastest drifter at
61--215~m/day, or 1.8--6.4~km per month, with PCO and GERO in between.
Passive collision safety and low station-keeping cost are therefore in
direct tension: the geometries that are safe to fly untended are the ones
that most need active control.

\begin{figure}[t]
  \centering
  \genfig{fig3_drift_rates}{0.95}{\texttt{python scripts/run\_drift.py}}
  \caption{Uncontrolled drift rate vs chief inclination and formation size.}
  \label{fig:drift}
\end{figure}
\begin{figure}[t]
  \centering
  \genfig{fig4_drift_contour}{0.55}{\texttt{python scripts/run\_drift.py}}
  \caption{Drift rate over $(\delta a, \delta i_x)$ offsets.}
  \label{fig:contour}
\end{figure}

\subsection{Perfect-state controller baselines}

With perfect navigation (Fig.~\ref{fig:perfect}, numerical truth), the
three controllers trace the expected cost/precision frontier. The
closed-form impulsive scheme is loose and cheap, at roughly 37~m/s/yr for
about 45~m RMS formation-keeping error; MPC occupies the middle at roughly
85~m/s/yr for about 16~m; and LQR buys meter-level control at steeply rising
cost, reaching $\sim$1~m RMS. Below about 0.9~m the LQR locus folds back on
itself, where tighter weighting stops improving accuracy and only burns
fuel, indicating a control-authority floor set by the saturation bound and
the 60-s actuation cadence. These curves are the reference against which all
estimation-in-the-loop cost increases are measured.

\begin{figure}[t]
  \centering
  \genfig{fig5_perfect_pareto}{0.55}{\texttt{python scripts/run\_perfect\_state.py}}
  \caption{Perfect-navigation $\Delta V$/yr vs RMS formation-keeping error
  (numerical truth).}
  \label{fig:perfect}
\end{figure}

\subsection{Estimation in the loop: the reliability map}

Running real filters produces a result that covariance-floor analyses cannot
generate: a large fraction of closed-loop trials fail outright. Across the
full tier, 22.1\% of LQR trials, 19.7\% of MPC trials, and 16.0\% of
impulsive trials diverge. Failures are not uniformly distributed. They
concentrate in angles-only navigation, where 52\% of S--S-filtered, 44\% of
KGD-filtered, and 23\% of GA-filtered trials diverge, and at specific
controller, sensor and sigma combinations the failure is total: 21 of 329
grid points have zero surviving seeds out of 100, concentrated in LQR and
MPC under angles-only navigation. At specific RF sigmas, MPC divergence
exceeds 90/100. These infeasible points are marked explicitly in every
figure ($\times$ symbols in Fig.~\ref{fig:headline}) rather than omitted,
since the location of the cliff is itself a design input. The practical
reading is that angles-only relative navigation, adequate in open-loop
estimation studies, frequently cannot close a 30-day formation-keeping loop
at 1-km scale with these controllers, which echoes flight experience that
angles-only operations require careful maneuver design
\cite{damico2013argon,gaias2014anglesonly}.

\begin{figure}[t]
  \centering
  \genfig{fig6_dv_vs_sigma}{0.95}{\texttt{python -m eilj2.figures.fig6\_dv\_vs\_sigma}}
  \caption{Closed-loop $\Delta V$/yr vs navigation accuracy (E/I-safe
  geometry, KGD filter, open-loop-calibrated $Q$). Shaded bands are
  interquartile ranges over 100 seeds; $\times$ marks sigma points at which
  all 100 seeds diverged. Note the \emph{falling} trend in the CDGPS panel,
  the signature of the filter-consistency problem analyzed in
  Sec.~\ref{sec:consistency}.}
  \label{fig:headline}
\end{figure}

\subsection{Filter consistency governs the apparent trade}
\label{sec:consistency}

Figure~\ref{fig:headline} contains a physically implausible result. Median
delta-V falls monotonically as CDGPS accuracy degrades from 5~mm to 10~cm
(for LQR/KGD, 85~$\to$~12~m/s/yr, Spearman $\rho = -1.00$), and
formation-keeping error improves at the same time, so that better sensors
appear to cost both fuel and accuracy. The explanation lies in the filter
rather than in the physics. At $\sigma = 5$~mm the KGD-filtered EKF has a
median NIS of $1.3\times10^4$ across CDGPS cells, against an expectation of
$\dim z = 3$, so the innovation covariance is understated by more than three
orders of magnitude. The filter over-weights its prediction, and the
controller chases a confidently wrong estimate, burning fuel in proportion
to the large innovation content of accurate measurements. As $\sigma$ grows,
the measurement weight shrinks toward what the broken covariance can
support and the pathology relaxes, so NIS falls toward consistency and the
apparent cost falls with it. The effect is ordered by process noise: at the
same 5-mm point the S--S filter ($q = 5.5\times10^{-10}$) has NIS 4.2, which
is consistent, while GA ($8.0\times10^{-13}$) reads 134 and KGD
($1.0\times10^{-14}$) reads $1.3\times10^4$. The open-loop calibration of
Section~\ref{sec:nav} gave the best dynamics model the least protection.

Two experiments pin down the mechanism. First, disabling control entirely
makes consistency 15$\times$ worse: in the LQR/along-track cell, whose
median NIS is $8.6\times10^3$ with control active, it rises to
$1.2\times10^5$ with the controller inhibited for the whole run. The
mismatch is therefore the STM's secular error (Fig.~\ref{fig:validation}),
which grows with unconstrained drift, and the control loop partially masks
it by holding the state near the reference; control feed-through is
exonerated. Second, re-tuning $q$ by closed-loop covariance matching,
inflating it until time-averaged NIS matches $\dim z$
(\texttt{scripts/tune\_q\_closedloop.py}), requires
$q_{\mathrm{KGD}}: 1.0\times10^{-14} \to 3.2\times10^{-10}$, a factor of
$3\times10^4$, which lands within a factor of two of the S--S value. The
required covariance is thus nearly common-mode across dynamics models of
vastly different fidelity, because in closed loop the drift that
distinguishes them is suppressed and what remains is the shared secular
character of the residual.

Figure~\ref{fig:qtuning} shows the consequence at $N = 100$. With the
NIS-matched filter, consistency holds across the entire sweep, with median
NIS between 3.1 and 5.9 against $\dim z = 3$, and the trade un-inverts. LQR
delta-V rises monotonically with sigma ($137 \to 162$~m/s/yr over
5~mm--10~cm, $\rho = +1.00$) and so does its formation-keeping error, which
is the physically expected behavior. MPC responds differently: its delta-V
is nearly flat while its RMS error rises monotonically, so MPC trades
accuracy rather than fuel for navigation quality, a controller-level
distinction that a single-metric study cannot reveal. Consistency carries a
cost. At the 5-mm point the consistent MPC loop holds 10.2~m RMS at
51~m/s/yr where the inconsistent one held 4.8~m at 19.6~m/s/yr, because
inflated process noise raises the filter gain and feeds measurement noise
into the control. Finally, even the matched filter's NEES remains between 28
and 1460, falling as $\sigma$ grows: the covariance is consistent in the
measured subspace but still optimistic in the unobserved directions, as
expected when a secular error is modeled as white noise. No white-$Q$ EKF
can be simultaneously consistent and well-performing against a secular model
error, and the principled remedy is to estimate the bias with a
Schmidt--Kalman or bias-augmented formulation, which we leave as follow-on
work. At the loosest CDGPS point ($\sigma = 1$~m) both corrected curves roll
off as the filter enters a measurement-rejection regime and the loop coasts
on the prediction, so the knee analysis below reports trends over the
flight-relevant 5~mm--10~cm range.

\begin{figure}[t]
  \centering
  \genfig{fig9_qtuning}{0.95}{\texttt{python -m eilj2.figures.fig9\_qtuning}}
  \caption{Process-noise calibration governs the apparent trade (KGD
  filter, CDGPS, along-track geometry, $N = 100$). Left: median NIS vs
  sigma for the open-loop-calibrated and NIS-matched $q$; dashed line marks
  consistency. Right: the resulting delta-V curves with interquartile
  bands and Spearman trend coefficients.}
  \label{fig:qtuning}
\end{figure}

\subsection{Knee summary}

Table~\ref{tab:knees} summarizes the robust knee metric of
Section~\ref{sec:kneemetric} over the full tier. The per-filter status
distribution makes the consistency contamination visible: under the
open-loop-calibrated $Q$, only 1 of 27 KGD cells shows the physically
expected rising trend, against 11 of 27 for GA and 6 of 27 for S--S, so the
better the dynamics model, the more thoroughly its trade curve is corrupted.
Where knees are measurable they are directly usable. With the consistent
S--S filter, LQR's CDGPS knee sits at 3--10~cm ($1\sigma$, threshold
metric), so centimeter-class relative GPS saturates the fuel benefit and
sub-centimeter accuracy buys nothing; under the NIS-matched KGD filter the
corrected LQR curve places the 10\%-rise point in the same decade. For
angles-only navigation the knees that exist sit at 10--50~arcsec, though the
reliability map above bounds their usefulness. Of 81 cells, 60 have
threshold-stable verdicts; the remainder are reported with their instability
flagged rather than forced to a point estimate.

\begin{table}[t]
  \centering
  \caption{Knee-metric status over the full tier (median $\Delta V$ trend
  vs $\sigma$ per cell; Sec.~\ref{sec:kneemetric}). ``Rising'' is the
  physically expected behavior for a consistent loop.}
  \label{tab:knees}
  \begin{tabular}{lccc}
    \toprule
    Filter model ($q$, open-loop cal.) & rising & flat & falling \\
    \midrule
    KGD ($1.0\times10^{-14}$) & 1 & 12 & 14 \\
    Gim--Alfriend ($8.0\times10^{-13}$) & 11 & 7 & 9 \\
    Schweighart--Sedwick ($5.5\times10^{-10}$) & 6 & 13 & 8 \\
    \bottomrule
  \end{tabular}
\end{table}

\subsection{Fidelity in the loop}

The campaign's filter-model dimension is itself the STM-fidelity ablation,
and its answer runs against expectation. Open loop, KGD beats S--S by a
factor of roughly 500 (Fig.~\ref{fig:validation}). In the closed loop, with
each filter carrying its conventionally calibrated $Q$, the crude S--S
filter is the only one that is statistically consistent without further
tuning, it produces the only systematically usable trade curves, and,
because closed-loop covariance demand is common-mode, it gives up far less
than the validation gap suggests. KGD retains the best point estimates, with
RMS position error of 3.8~m against 13--15~m for the others under CDGPS, so
the high-fidelity model is the right choice when its covariance is matched
in closed loop; fidelity spent on the mean while the covariance is starved
gives the worst of both. For a mission designer, the question of which model
the filter should carry is therefore inseparable from how its process noise
will be calibrated, which is the central practical lesson of this study. The
Pareto surface over navigation accuracy and control tightness
(Fig.~\ref{fig:surface}) completes the picture for the consistent-filter
configuration.

\begin{figure}[t]
  \centering
  \genfig{fig7_pareto_surface}{0.95}{\texttt{python -m eilj2.figures.fig7\_pareto\_surface}}
  \caption{Pareto surface over (nav accuracy, control precision,
  $\Delta V$).}
  \label{fig:surface}
\end{figure}
